% Трёхфазная нулевая схема (три диода) — CircuitTikZ, в стиле bridge_circuitikz.tex
% Три одинаковые вертикальные ветви (обмотка + диод на одной высоте),
% точки-узлы только в T-соединениях
\documentclass[border=4pt]{standalone}
\usepackage{amsmath}
\usepackage{fontspec}
\setmainfont{Calibri}
\usepackage[european, RPvoltages]{circuitikz}
\definecolor{ink}{RGB}{31,42,51}
\definecolor{grp}{RGB}{46,110,140}
\ctikzset{bipoles/length=1.0cm}
\ctikzset{bipoles/thickness=0.9}
\ctikzset{diodes/fill=ink}      % залитые диоды (ГОСТ)
\ctikzset{inductor=american}    % витки-полуокружности

\begin{document}
\begin{tikzpicture}[line width=1.0pt, color=ink, font=\small]
  % --- три одинаковые вертикальные ветви: обмотка + диод на одной высоте
  \foreach \x/\name/\ph in {0/VD1/a, 2.2/VD2/b, 4.4/VD3/c} {
    \draw (\x,-1.8) to[L] (\x,0.6);
    \draw (\x,0.6) -- (\x,2.4) to[D] (\x,3.0) -- (\x,4.6);
    \node[anchor=west] at (\x+0.15,-0.6) {$\ph$};
    \node[anchor=west] at (\x+0.35,2.7) {\name};
  }
  % --- нижняя шина с нулевым выводом N (угол a — без точки)
  \draw (0,-1.8) -- (10.0,-1.8);
  \node[anchor=east] at (-0.3,-1.8) {$N$};
  \node[circ] at (2.2,-1.8) {};
  \node[circ] at (4.4,-1.8) {};
  % --- верхняя шина + (угол VD1 — без точки)
  \draw (0,4.6) -- (10.0,4.6);
  \node[circ] at (2.2,4.6) {};
  \node[circ] at (4.4,4.6) {};
  % --- нагрузка
  \draw (10.0,4.6) -- (10.0,3.0) to[R] (10.0,1.0) -- (10.0,-1.8);
  % --- подписи
  \node[anchor=east] at (9.55,2.0) {$u_d$};
  \node[anchor=west] at (10.45,2.0) {$R_{\text{н}}$};
  \node[anchor=south west] at (9.95,4.62) {$+$};
  \node[anchor=north west] at (9.95,-1.62) {$-$};
\end{tikzpicture}
\end{document}
